\documentclass[12pt]{article}
\usepackage{sbc-template}
\usepackage[T1]{fontenc}
\usepackage[utf8]{inputenc}
\usepackage[brazil]{babel}
\usepackage{graphicx,url,amsmath,amssymb,booktabs,tabularx}
\sloppy
\title{Inteligência Artificial no Ensino dos Fundamentos de Programação Competitiva:\\Recuperação, Validação e Limitações sob Orçamento Zero de API}
\author{Autoria a preencher\inst{1}}
\address{Instituição e departamento a preencher\\Cidade -- Estado -- Brasil\email{email@instituicao.br}}
\begin{document}
\maketitle
\begin{abstract}
This paper presents AutoTutor as an exploratory prototype for learning competitive programming foundations. It combines language models, retrieval-augmented generation, persistent learning artifacts, and executable verification. We map implemented mechanisms to AutoCode and OI-Assistant, with Codigen as an additional reference. A retrospective analysis of four foundational tasks exposes interpretation, output-format and truncation failures in a lightweight local model. A separate API pilot verifies connectivity and validation infrastructure but does not evaluate beginner learning outcomes. We propose a controlled comparison of retrieval and generation strategies using Gemini 3.1 Flash-Lite within the user's declared free quota. The contribution is an auditable architecture and evaluation protocol, without claims of educational effectiveness, general model superiority or fine-tuning.
\end{abstract}
\begin{resumo}
Este artigo apresenta o AutoTutor como protótipo exploratório para aprendizagem dos fundamentos de programação competitiva. A aplicação combina modelos de linguagem, geração aumentada por recuperação, persistência de materiais e verificação executável. Mapeiam-se mecanismos implementados aos trabalhos AutoCode e OI-Assistant, com Codigen como referência adicional. Uma análise retrospectiva de quatro tarefas básicas expõe falhas de interpretação, formato e truncamento de um modelo local leve. Um piloto separado de API verifica conectividade e infraestrutura de validação, sem avaliar aprendizagem de iniciantes. Propõe-se comparação controlada de estratégias com Gemini 3.1 Flash-Lite na cota gratuita declarada pelo usuário. A contribuição é uma arquitetura auditável e um protocolo de avaliação, sem alegações de eficácia educacional, superioridade geral ou ajuste de pesos.
\end{resumo}

\section{Introdução e delimitação}
A formação inicial em programação competitiva envolve interpretar enunciados, percorrer vetores, distinguir casos, estimar complexidade e relacionar uma solução ao resultado esperado. Um iniciante pode compreender uma explicação aparentemente convincente e ainda construir um índice incorreto ou confundir a ordem de visitação de um grafo. Assistentes de IA podem oferecer exemplos e dicas, mas a fluência do texto não assegura a correção do procedimento.

O AutoTutor reúne estudo, biblioteca, problemas, juiz local e acompanhamento em uma aplicação Python/Streamlit com persistência SQLite. O presente recorte prioriza fundamentos e simulação Ad Hoc, vetores e contagem, ordenação, somas de prefixos, busca binária, guloso introdutório, BFS, DFS e programação dinâmica elementar. Árvores de segmentos, propagação preguiçosa e Block-Cut Tree ficam fora do conjunto principal de avaliação. A existência de recursos avançados na aplicação não define o público nem a contribuição deste trabalho.

A pergunta orientadora é: \emph{como recuperação de referências e validação independente podem apoiar a produção de explicações e exercícios introdutórios, considerando limitações de modelos locais e uma API sem orçamento financeiro?} As contribuições são: (i) mapear os mecanismos efetivamente implementados e sua relação com a literatura; (ii) analisar falhas locais de forma rastreável; e (iii) estabelecer um protocolo viável de avaliação com orçamento zero. Este texto relata engenharia e evidências exploratórias. Não apresenta um experimento concluído com estudantes.

\section{Fundamentação de inteligência artificial}
\subsection{Modelos de linguagem, inferência e recursos}
Um modelo autorregressivo descreve a distribuição de uma resposta $y$ condicionada ao contexto $x$ por
\begin{equation}
p_\theta(y\mid x)=\prod_{t=1}^{T}p_\theta(y_t\mid x,y_{<t}).
\end{equation}
Os parâmetros $\theta$ decorrem do treinamento. O AutoTutor realiza inferência com modelos existentes, sem alterar seus pesos. Salvar um estudo ou adicionar documentos ao banco não constitui treinamento, aprendizagem contínua ou \emph{fine-tuning}. O Transformer fundamenta essa família de modelos pelo mecanismo de atenção, que relaciona representações de posições do contexto \cite{vaswani2017}. A aplicação utiliza o serviço de inferência e não implementa a arquitetura interna do provedor.

O Llama 3.2 1B permanece como opção local por restrição de recursos. Um modelo compacto pode operar na máquina disponível, mas sua capacidade pedagógica precisa ser medida. O custo de memória inclui pesos, contexto, estruturas intermediárias e processo de execução. Assim, tamanho do arquivo e número de parâmetros não descrevem isoladamente a RAM necessária. A configuração local limita threads e solicita descarregamento após a resposta. Esses controles afetam a operação, sem certificar correção.

O acesso por API transfere a inferência para outra infraestrutura. Pode reduzir pressão de memória no computador, mas acrescenta dependência de rede, cota, indisponibilidade e políticas do provedor. Não é possível deduzir a superioridade de um modelo apenas de sua localização. Temperatura reduzida e saída JSON também não garantem verdade: regulam variabilidade e forma, respectivamente.

\subsection{RAG, conhecimento externo e proveniência}
RAG condiciona a geração a evidências recuperadas de um acervo \cite{lewis2020}. No protótipo, o fluxo pode ser descrito como
\begin{equation}
 D_k=\operatorname{Retrieve}(q,\mathcal{D}),\qquad
 y\sim p_\theta(\cdot\mid q,D_k,s),
\end{equation}
em que $q$ representa o pedido, $\mathcal{D}$ o corpus, $D_k$ os trechos selecionados e $s$ as instruções. O recuperador e o gerador não são treinados conjuntamente. O armazenamento expande a memória externa disponível na inferência.

Há falhas distintas de cobertura, recuperação e utilização. Um documento pode faltar no corpus, existir sem aparecer entre os primeiros resultados ou ser recuperado e interpretado incorretamente. URLs no rodapé não demonstram que o conteúdo foi utilizado. Por isso, o sistema diferencia texto integral externo, material autoral e referência sem texto, registrando origem e trechos recuperados. Materiais gerados por IA permanecem disponíveis para estudo posterior, mas não substituem automaticamente fontes externas na recuperação experimental.

BM25 combina evidência lexical, saturação de frequência e normalização de comprimento \cite{robertson2009}. É econômico e auditável, mas pode não recuperar paráfrases com vocabulário distinto. Embeddings, motivados por representações como Sentence-BERT \cite{reimers2019}, aproximam textos por vetores. A similaridade cosseno $u^\top v/(\|u\|\|v\|)$ mede proximidade representacional, sem provar equivalência algorítmica.

O experimento implementa recuperação lexical, densa, híbrida, reordenação leve, diversificação e expansão por competências. A fusão RRF soma $1/(60+\operatorname{pos}(d))$ sobre listas, com posições iniciando em um nesta versão \cite{cormack2009}. O re-ranking atual usa evidência lexical e metadados, sem um modelo neural de relevância. A diversificação usa uma aproximação por sobreposição de palavras. Nenhuma dessas opções demonstrou superioridade no recorte introdutório. O objetivo inicial é avaliar se melhoram evidências úteis, antes de adotar infraestrutura mais complexa.

\subsection{Competências e evidências do estudante}
Uma competência representa um objetivo observável, como construir prefixos sem erro de índice ou usar uma fila para calcular distâncias sem pesos. Um grafo editorial relaciona 18 competências a pré-requisitos recomendados e contratos. Essas relações orientam consulta e especificação, sem estabelecer uma única ordem obrigatória de aprendizagem. A revisão por professores permanece pendente.

O protótipo BKT mantém uma probabilidade latente de conhecimento com parâmetros de conhecimento inicial, aprendizagem, acerto por acaso e erro apesar do domínio \cite{corbett1994}. Para resposta correta, uma atualização simplificada é
\begin{equation}
 P(L\mid C)=\frac{P(L)(1-P(S))}{P(L)(1-P(S))+(1-P(L))P(G)},
\end{equation}
seguida da transição $P(L')=P(L\mid C)+(1-P(L\mid C))P(T)$. Aqui os parâmetros são hipotéticos e não foram ajustados a estudantes. Repetir tentativas do mesmo item não multiplica evidência independente. Menos de três itens distintos permanece insuficiente, e esse mínimo é uma decisão editorial, sem validade diagnóstica demonstrada.

Objetivos graduais, dicas e feedback motivam o desenho pedagógico, mas são diferentes de uma técnica de IA e de uma medida de aprendizagem. Um veredito WA pode resultar de índice, interpretação ou implementação. O modelo não deve convertê-lo automaticamente em ausência de conhecimento. Contagens e taxas devem ser calculadas por software e fornecidas ao modelo como evidências; gerar uma interpretação não autoriza inventar estatísticas.

\section{Trabalhos relacionados e mecanismos aproveitados}
O AutoCode propõe ciclos de geração e validação de problemas e testes \cite{zhou2026}. O AutoTutor aproveita a ideia de confrontar soluções e procurar inconsistências, mas restringe o experimento a contratos conhecidos de entrada e saída. São implementados compilação isolada, exemplos e comparação com oráculos e referências confiáveis. Faltam um benchmark de vereditos oficiais, revisão de originalidade e avaliação equivalente à do artigo. A consistência próxima de 99\% relatada no AutoCode pertence à avaliação daquele trabalho e não corresponde à eficácia educacional nem à acurácia deste protótipo.

O OI-Assistant integra recuperação de problemas semelhantes e aprendizagem interativa \cite{su2025}. O AutoTutor distingue recomendação de problemas e recuperação de materiais: a primeira utiliza representações lexicais e conceituais; a segunda usa a biblioteca externa e contexto com proveniência. Recuperação por competências é uma extensão experimental local, sem alegar reprodução integral do OI-Assistant.

Codigen relaciona autoria e personalização com modelos de linguagem \cite{codigen2025}. É uma referência adicional para a edição do pedido, profundidade do estudo e apresentação do conteúdo. A aplicação não reproduz sua interface de autoria visual nem seus resultados. A contribuição do AutoTutor é integrar e tornar observáveis mecanismos, não apresentar um novo modelo fundacional.

\section{Arquitetura, estado atual e melhorias propostas}
A aplicação persiste usuários, turmas, materiais, trilhas e submissões em banco local. Páginas próprias permitem retomar estudos, anotações e referências. O acompanhamento do professor e a proposta de trios são recursos da plataforma, fora das variáveis centrais desta avaliação. Sua existência não demonstra que a análise automática esteja correta.

O corpus contém 5551 registros: 302 textos externos, 19 conteúdos autorais e 5230 referências. Há 6837 chunks no banco, incluindo diferentes categorias. Portanto, a quantidade de referências não deve ser apresentada como quantidade de artigos integralmente lidos. Os textos externos incluem 164 páginas de CP-Algorithms e 138 de USACO Guide, com licenças e origem. Não se afirma importação integral de todos os sites solicitados \cite{cpalgorithms2026,usaco2026}.

Uma auditoria adicional de nove consultas básicas no BM25 identificou uma falha de vocabulário: ``Vetores e contagem'' não retornava texto externo. A inclusão de equivalências português--inglês passou a recuperar candidatos, sem demonstrar relevância pedagógica. Outras consultas retornaram tópicos como matroides, 0--1 BFS e DP de perfil entre os primeiros resultados. Essa inspeção de fontes evidencia necessidade de curadoria por nível; não substitui julgamentos completos de relevância nem autoriza calcular precisão e recall.

O banco local possui 39 contratos executáveis. Nesta revisão, nove objetivos básicos receberam associação explícita a contratos; isso não significa que todos os 39 sejam introdutórios. Soma de prefixos passou a ser o padrão do piloto, e o painel experimental oculta tópicos avançados até seleção explícita. Um teste offline do contrato de prefixos passou pela validação, utilizando solução conhecida. Ele verifica integração, não qualidade de geração por modelo.

O caminho experimental segue: objetivo e preferências; contrato e competências; recuperação; geração estruturada; schema; compilação e testes; armazenamento com fontes e validação. O juiz usa namespaces Linux, rede ausente, filesystem restrito e limites de tempo, memória, saída e processos. A ausência de isolamento bloqueia execução. Trata-se de protótipo local, sem auditoria para exposição pública multiusuário.

As próximas melhorias prioritárias são curar fontes introdutórias por objetivo, reduzir recuperação de conteúdo avançado irrelevante, verificar exemplos com ferramentas determinísticas e avaliar dicas por rubrica humana. O índice semântico e o grafo devem justificar seu custo por resultados mensuráveis. Cache local de materiais já validados pode economizar chamadas, preservando autoria, modelo e versão; o cache de respostas por chave de conteúdo permanece proposta, não implementação atual.

\section{Evidências observadas e limites dos modelos locais}
\subsection{Reanálise de tarefas básicas com Llama 1B}
Uma avaliação anterior registrou cinco tarefas sintéticas com Llama 3.2 1B e limite de 256 tokens de saída. Para o novo recorte, considera-se apenas BFS, DFS, DP e prefixos; o item avançado fica excluído. As respostas recebiam referências externas. A Tabela~\ref{tab:local} distingue modos de falha e evita reduzir todos a erro aritmético.

\begin{table}[ht]
\centering\small
\caption{Quatro tarefas básicas da avaliação anterior. Tempos observados não isolam carga concorrente.}\label{tab:local}
\begin{tabularx}{\textwidth}{l r X}
\toprule Tema & Tempo (s) & Evidência registrada \\
\midrule
BFS & 24,82 & Esperava-se a sequência de distâncias $[0,1,1,2]$. A saída forneceu pares diferentes e confundiu a descrição da busca. \\
DFS & 17,85 & Esperava-se 4; a resposta forneceu resultado 1 e rótulos de fontes em formato inválido. \\
DP & 34,04 & A saída truncou durante uma enumeração, impedindo leitura do JSON. Não isola falha de raciocínio. \\
Prefixos & 26,38 & O valor 5 aparece, mas a justificativa apresenta fórmula/resultado contraditórios e o campo esperado não é respeitado. \\
\bottomrule
\end{tabularx}
\end{table}

O verificador original não aprovou os quatro itens. Isso mede conformidade com o contrato e consistência da resposta nessas condições, não uma acurácia geral de 0\% em programação. O limite curto pode explicar o truncamento de DP. Para prefixos, um número coincidente não corrige a justificativa nem o schema. A ausência de repetições e de controle de carga impede inferência estatística sobre o modelo.

O Qwen2.5 3B também foi testado anteriormente em perfis sintéticos. Houve geração de material que excedeu 300 segundos, inversão de taxas e justificativas com informações não fornecidas. A configuração e as tarefas diferem das do Llama. Portanto, não se constitui ranking entre famílias. O caso evidencia que JSON válido e IDs sem repetição não asseguram fidelidade a dados. Um teste de conectividade do Llama respondeu ``OK'' em 3,81 segundos com tamanho de modelo reportado de 1466852966 bytes; esse valor não é medição de RAM total da aplicação.

\subsection{O que o piloto de API permite afirmar}
O adaptador Gemini 3.1 Flash-Lite recebeu resposta em quatro chamadas reais anteriores, incluindo conexão, somando 5744 tokens retornados pelo provedor. O piloto de geração utilizou um contrato avançado e não compõe a avaliação de fundamentos. Seus artefatos ficam preservados como evidência de infraestrutura. Não foi executada nesta revisão uma comparação pareada Gemini--Llama em problemas básicos. Consequentemente, a API está operacional, mas sua vantagem pedagógica no público-alvo permanece hipótese.

\section{Gratuidade e protocolo de avaliação proposto}
A documentação consultada em 5 de novembro de 2026 lista entrada e saída do Gemini 3.1 Flash-Lite sem cobrança na modalidade gratuita. Cotas dependem do projeto e devem ser consultadas no AI Studio \cite{googlefree2026}. A aplicação não identifica o plano de faturamento pelo texto de uma resposta. O usuário declarou utilizar cota gratuita, e não há autorização para despesa. O adaptador não habilita faturamento nem usa fallback pago.

O limite de seis tentativas por dia UTC é uma proteção local do projeto, não a cota publicada pelo Google. Também há teto de tokens reservados, entrada limitada, intervalo mínimo e interrupção após erros de cota. Não se deve confundir o reinício UTC do ledger com o reinício de cota do provedor. Falhas continuam no denominador e consomem reservas. Grounding por busca e cache gerenciado não são necessários: o RAG usa SQLite e embeddings locais, sem depender desses serviços.

Os termos de serviços gratuitos permitem uso de entradas e saídas para melhoria de produtos e revisão. Assim, esta fase utiliza somente dados públicos e sintéticos, sem nomes, códigos privados ou históricos identificáveis de estudantes \cite{googleterms2026}. Uma implantação educacional real exige revisar os termos aplicáveis e o tratamento de dados antes de enviar informações de participantes.

Propõe-se um primeiro piloto com quatro temas (simulação, prefixos, BFS e DP introdutória) e três condições: contrato sem recuperação, contrato com BM25 e contrato com BM25 mais especificação pedagógica. São 12 chamadas sem reparo. Mantendo o teto de seis, seriam necessários pelo menos dois dias com a capacidade diária integral disponível, condicionados às cotas reais. Repetir cada célula três vezes demanda 36 chamadas e pelo menos seis dias nessas condições. Esses números são planejamento de chamadas, não promessa de disponibilidade gratuita ou resultado executado.

O piloto compara pacotes de contexto. Para isolar a contribuição da especificação, uma ablação posterior manterá o contexto recuperado idêntico. Ordem aleatória, tarefas pareadas, mesmo modelo, temperatura, orçamento de saída e critérios reduzem confundimento. Um comparativo local/API adicional deve padronizar prompts, schema, limite de saída e timeout, registrar versões e separar falhas de formato, raciocínio, truncamento e execução. A configuração local antiga não oferece esse controle.

As métricas técnicas incluem schema, compilação, exemplos e testes independentes. A rubrica humana cobre clareza, ambiguidade, necessidade da competência, coerência das restrições, progressão e dicas. Precision, Recall, MRR e nDCG exigem julgamentos de relevância revisados; não foram calculados sem rótulos. Dois revisores independentes devem avaliar saídas sem conhecer a condição. Taxas devem indicar denominador e incerteza. Ganho de aprendizagem exige estudo separado, com instrumentos e participantes adequados.

\section{Ameaças à validade e conclusão}
As tarefas locais são poucas e sintéticas, com limite de saída curto e sem repetições. Dados de API disponíveis referem-se a outro contrato. O corpus tem cobertura desigual e inclui materiais avançados. O grafo e o BKT são editoriais, sem validação especializada ou calibração. Testes finitos não provam correção nem certificam qualidade pedagógica, e contratos conhecidos limitam originalidade algorítmica.

O AutoTutor constitui uma base funcional para investigar apoio aos fundamentos: fontes com proveniência, geração configurável, persistência e validação executável estão disponíveis. A análise local mostra por que modelos leves precisam de verificações, inclusive em tarefas simples. A cota gratuita permite planejar experimentos pequenos e progressivos, sem garantir escala de turma ou reproduzir resultados de sistemas de maior capacidade. A próxima contribuição científica depende de executar o protocolo básico e revisar suas saídas, preservando a distinção entre implementação, evidência técnica e aprendizagem.

\bibliographystyle{sbc}
\bibliography{referencias}
\end{document}
